\documentclass[12pt,a4paper,openright,dvipsnames]{mwrep}

%%% libraries

% ~egreg | https://tex.stackexchange.com/a/95842

\usepackage{array, tabularx}

\newenvironment{conditions}
  {\par\vspace{\abovedisplayskip}\noindent
   \begin{tabular}{>{$}l<{$} @{} >{${}}c<{{}$} @{} l}}
  {\end{tabular}\par\vspace{\belowdisplayskip}}

\newenvironment{conditions*}
  {\par\vspace{\abovedisplayskip}\noindent
   \tabularx{\columnwidth}{>{$}l<{$} @{}>{${}}c<{{}$}@{} >{\raggedright\arraybackslash}X}}
  {\endtabularx\par\vspace{\belowdisplayskip}} % where: conditions
% ~ignasi | https://tex.stackexchange.com/a/140581

\usepackage{tikz}
\usetikzlibrary{matrix,calc}

%isolated term
%#1 - Optional. Space between node and grouping line. Default=0
%#2 - node
%#3 - filling color
\newcommand{\implicantsol}[3][0]{
    \draw[rounded corners=3pt, fill=#3, opacity=0.3] ($(#2.north west)+(135:#1)$) rectangle ($(#2.south east)+(-45:#1)$);
}


%internal group
%#1 - Optional. Space between node and grouping line. Default=0
%#2 - top left node
%#3 - bottom right node
%#4 - filling color
\newcommand{\implicant}[4][0]{
    \draw[rounded corners=3pt, fill=#4, opacity=0.3] ($(#2.north west)+(135:#1)$) rectangle ($(#3.south east)+(-45:#1)$);
}

%group lateral borders
%#1 - Optional. Space between node and grouping line. Default=0
%#2 - top left node
%#3 - bottom right node
%#4 - filling color
\newcommand{\implicantcostats}[4][0]{
    \draw[rounded corners=3pt, fill=#4, opacity=0.3] ($(rf.east |- #2.north)+(90:#1)$)-| ($(#2.east)+(0:#1)$) |- ($(rf.east |- #3.south)+(-90:#1)$);
    \draw[rounded corners=3pt, fill=#4, opacity=0.3] ($(cf.west |- #2.north)+(90:#1)$) -| ($(#3.west)+(180:#1)$) |- ($(cf.west |- #3.south)+(-90:#1)$);
}

%group top-bottom borders
%#1 - Optional. Space between node and grouping line. Default=0
%#2 - top left node
%#3 - bottom right node
%#4 - filling color
\newcommand{\implicantdaltbaix}[4][0]{
    \draw[rounded corners=3pt, fill=#4, opacity=0.3] ($(cf.south -| #2.west)+(180:#1)$) |- ($(#2.south)+(-90:#1)$) -| ($(cf.south -| #3.east)+(0:#1)$);
    \draw[rounded corners=3pt, fill=#4, opacity=0.3] ($(rf.north -| #2.west)+(180:#1)$) |- ($(#3.north)+(90:#1)$) -| ($(rf.north -| #3.east)+(0:#1)$);
}

%group corners
%#1 - Optional. Space between node and grouping line. Default=0
%#2 - filling color
\newcommand{\implicantcantons}[2][0]{
    \draw[rounded corners=3pt, opacity=.3] ($(rf.east |- 0.south)+(-90:#1)$) -| ($(0.east |- cf.south)+(0:#1)$);
    \draw[rounded corners=3pt, opacity=.3] ($(rf.east |- 8.north)+(90:#1)$) -| ($(8.east |- rf.north)+(0:#1)$);
    \draw[rounded corners=3pt, opacity=.3] ($(cf.west |- 2.south)+(-90:#1)$) -| ($(2.west |- cf.south)+(180:#1)$);
    \draw[rounded corners=3pt, opacity=.3] ($(cf.west |- 10.north)+(90:#1)$) -| ($(10.west |- rf.north)+(180:#1)$);
    \fill[rounded corners=3pt, fill=#2, opacity=.3] ($(rf.east |- 0.south)+(-90:#1)$) -|  ($(0.east |- cf.south)+(0:#1)$) [sharp corners] ($(rf.east |- 0.south)+(-90:#1)$) |-  ($(0.east |- cf.south)+(0:#1)$) ;
    \fill[rounded corners=3pt, fill=#2, opacity=.3] ($(rf.east |- 8.north)+(90:#1)$) -| ($(8.east |- rf.north)+(0:#1)$) [sharp corners] ($(rf.east |- 8.north)+(90:#1)$) |- ($(8.east |- rf.north)+(0:#1)$) ;
    \fill[rounded corners=3pt, fill=#2, opacity=.3] ($(cf.west |- 2.south)+(-90:#1)$) -| ($(2.west |- cf.south)+(180:#1)$) [sharp corners]($(cf.west |- 2.south)+(-90:#1)$) |- ($(2.west |- cf.south)+(180:#1)$) ;
    \fill[rounded corners=3pt, fill=#2, opacity=.3] ($(cf.west |- 10.north)+(90:#1)$) -| ($(10.west |- rf.north)+(180:#1)$) [sharp corners] ($(cf.west |- 10.north)+(90:#1)$) |- ($(10.west |- rf.north)+(180:#1)$) ;
}

%Empty Karnaugh map 4x4
\newenvironment{Karnaugh}[2]%
{
    \begin{tikzpicture}[baseline=(current bounding box.north),scale=0.8]
        \draw (0,0) grid (4,4);
        \draw (0,4) -- node [pos=1,above right,anchor=south west] {#1} node [pos=0.7,below left,anchor=north east] {#2} ++(135:1);
        %
        \matrix (mapa) [matrix of nodes,
            column sep={0.8cm,between origins},
            row sep={0.8cm,between origins},
            every node/.style={minimum size=0.3mm},
            anchor=8.center,
            ampersand replacement=\&] at (0.5,0.5)
        {
            \& |(c00)| 00         \& |(c01)| 01         \& |(c11)| 11         \& |(c10)| 10         \& |(cf)| \phantom{00}    \\
            |(r00)| 00             \& |(0)|  \phantom{0} \& |(1)|  \phantom{0} \& |(3)|  \phantom{0} \& |(2)|  \phantom{0} \& \\
            |(r01)| 01             \& |(4)|  \phantom{0} \& |(5)|  \phantom{0} \& |(7)|  \phantom{0} \& |(6)|  \phantom{0} \& \\
            |(r11)| 11             \& |(12)| \phantom{0} \& |(13)| \phantom{0} \& |(15)| \phantom{0} \& |(14)| \phantom{0} \& \\
            |(r10)| 10             \& |(8)|  \phantom{0} \& |(9)|  \phantom{0} \& |(11)| \phantom{0} \& |(10)| \phantom{0} \& \\
            |(rf) | \phantom{00}   \&                    \&                    \&                    \&                    \& \\
        };
        }%
        {
    \end{tikzpicture}
}

%Empty Karnaugh map 2x4
\newenvironment{Karnaughvuit}[2]%
{
    \begin{tikzpicture}[baseline=(current bounding box.north),scale=0.8]
        \draw (0,0) grid (4,2);
        \draw (0,2) -- node [pos=1,above right,anchor=south west] {#1} node [pos=0.7,below left,anchor=north east] {#2} ++(135:1);
        %
        \matrix (mapa) [matrix of nodes,
            column sep={0.8cm,between origins},
            row sep={0.8cm,between origins},
            every node/.style={minimum size=0.3mm},
            anchor=4.center,
            ampersand replacement=\&] at (0.5,0.5)
        {
            \& |(c00)| 00         \& |(c01)| 01         \& |(c11)| 11         \& |(c10)| 10         \& |(cf)| \phantom{00}   \\
            |(r00)| 0             \& |(0)|  \phantom{0} \& |(1)|  \phantom{0} \& |(3)|  \phantom{0} \& |(2)|  \phantom{0} \& \\
            |(r01)| 1             \& |(4)|  \phantom{0} \& |(5)|  \phantom{0} \& |(7)|  \phantom{0} \& |(6)|  \phantom{0} \& \\
            |(rf) | \phantom{00}  \&                    \&                    \&                    \&                    \& \\
        };
        }%
        {
    \end{tikzpicture}
}

%Empty Karnaugh map 2x2
\newenvironment{Karnaughquatre}[2]%
{
    \begin{tikzpicture}[baseline=(current bounding box.north),scale=0.8]
        \draw (0,0) grid (2,2);
        \draw (0,2) -- node [pos=0.7,above right,anchor=south west] {#1} node [pos=0.7,below left,anchor=north east] {#2} ++(135:1);
        %
        \matrix (mapa) [matrix of nodes,
            column sep={0.8cm,between origins},
            row sep={0.8cm,between origins},
            every node/.style={minimum size=0.3mm},
            anchor=2.center,
            ampersand replacement=\&] at (0.5,0.5)
        {
            \& |(c00)| 0          \& |(c01)| 1                    \\
            |(r00)| 0 \& |(0)|  \phantom{0} \& |(1)|  \phantom{0} \\
            |(r01)| 1 \& |(2)|  \phantom{0} \& |(3)|  \phantom{0} \\
        };
        }%
        {
    \end{tikzpicture}
}

%Defines 8 or 16 values (0,1,X)
\newcommand{\contingut}[1]{%
    \foreach \x [count=\xi from 0]  in {#1}
    \path (\xi) node {\x};
}

%Places 1 in listed positions
\newcommand{\minterms}[1]{%
    \foreach \x in {#1}
    \path (\x) node {1};
}

%Places 0 in listed positions
\newcommand{\maxterms}[1]{%
    \foreach \x in {#1}
    \path (\x) node {0};
}

%Places X in listed positions
\newcommand{\indeterminats}[1]{%
    \foreach \x in {#1}
    \path (\x) node {X};
} % karnaugh maps


%%% packages and basic settings

\usepackage{lmodern} % fixes poor font quality (use only if NOT using XeLaTeX or LuaLaTeX)
\usepackage[polish]{babel}
\usepackage[T1]{polski}
\usepackage[utf8]{inputenc}

\usepackage[a4paper,
            tmargin=2cm,
            bmargin=2cm,
            lmargin=2cm,
            rmargin=2cm,
            bindingoffset=0cm]{geometry}

% ToC and hyper links
\usepackage{tocloft} % controlling the typographic design of the Table of Contents
\usepackage{hyperref} % handle cross-referencing commands to produce hypertext links in the document

% maths
\usepackage{amsmath} % improves the information structure and printed output
\usepackage{amssymb} % provides an extended symbol collection
\usepackage{siunitx} % units \SI \si
% \usepackage{amsthm} % helps to define theorem-like structures

% code syntax
\usepackage{minted} % syntax highlighting using pygments (check docs for changing themes and styling)

% graphics and styling
\usepackage{graphicx} % graphics support
\usepackage{subfig}
\usepackage{float} % improves the interface for defining floating objects
\usepackage{setspace} % sets space between lines

% tables (check docs for every packages as each has a different task)
\usepackage{booktabs}
\usepackage{array}
\usepackage{makecell}
\usepackage{tabularx}

% coloring utility
\usepackage{xcolor}

% advanced graphic stuff
\usepackage{tikz}
\usetikzlibrary{automata, positioning, arrows} % automata diagrams

% black hyperlinks (change to blue for highlighting)
\hypersetup{
    colorlinks,
    citecolor=black,
    filecolor=black,
    linkcolor=black,
    urlcolor=black
}


%%% variables

\newcommand{\vtitle}{Title}
\newcommand{\vauthors}{
    Author One\\
    Author Two
}



%%% code starts here


% title page AGH
\begin{document}
\begin{titlepage}

    \centering
    \singlespacing

    \includegraphics[scale=0.5]{common/agh.png}

    \vspace*{0.4cm}

    \fontsize{12}{16}{\textbf{\MakeUppercase{Akademia Górniczo-Hutnicza}}}\\
    \fontsize{10}{12}{\textbf{\MakeUppercase{Im. Stanisława Staszica w Krakowie}}}

    \vspace*{0.4cm}

    \vspace{\stretch{0.3}}

    {\LARGE\bfseries \vtitle\\}
    \rule{3in}{0.4pt} \\
    \large\vauthors


    \vspace*{\stretch{2}}

    \centerline{\rmfamily\fontsize{14}{16}\selectfont Kraków, \today}

\end{titlepage}
\newpage

% ToC
\tableofcontents
\newpage


%
\chapter{Rozdział}

%%
\section{Sekcja 1}

%%%
\subsection{Syntax highlight}

\phantomsection\hypertarget{python-code}{}
\begin{minted}{python}
def getFirstPrimeNumbers(n):
    """Returns a list of the first n prime numbers"""
    primes = []
    i = 2
    while len(primes) < n:
        if isPrime(i):
            primes.append(i)
        i += 1
    return primes



def isPrime(n):
    """Returns True if n is prime, False otherwise"""
    for i in range(2, n):
        if n % i == 0:
            return False
    return True

if __name__ == "__main__":
    print(getFirstPrimeNumbers(10))
\end{minted}
\newpage

%%%
\subsection{Images}

\begin{figure}[H] % forces figure to this position (does not allow LaTeX to relocate for best fit)
    \centering
    \includegraphics[width=\textwidth]{example-image-a}
    \caption{example A image}
    \label{fig:example-a}
\end{figure}

\begin{figure}[H]
    \centering
    \subfloat[example subfig B]{%
        \includegraphics[width=0.5\textwidth]{example-image-b}%
    }
    \subfloat[example subfig C]{%
        \includegraphics[width=0.5\textwidth]{example-image-c}%
    }
    \caption{Subfigures B and C}
\end{figure}


%%%
\subsection{Phantom hyperlink}

Znalazłem interesujący \hyperlink{python-code}{kod} rozwiązujący nasz problem.


%%%
\subsection{Conditions}

Short conditions:
\begin{equation}
    \rho S \Delta x
    \frac{\partial^2 \Psi(x, t)}{\partial t^2}
    =
    S(\sigma (x + \Delta x) - \sigma(x))
    \text{,}
\end{equation}
gdzie:
\begin{conditions}
    S & - & powierzchnia przekroju pręta, \\
    \rho & - & gęstość materiału, \\
    \sigma & - & naprężenie, \\
    \Psi(x,t) & - & średnie przemieszczenie wycinka pręta.
\end{conditions}

\vspace*{0.6cm}

Long conditions:
\begin{equation}
    \rho S \Delta x
    \frac{\partial^2 \Psi(x, t)}{\partial t^2}
    =
    S(\sigma (x + \Delta x) - \sigma(x))
    \text{,}
\end{equation}
gdzie:
\begin{conditions*}
    S & - & powierzchnia przekroju pręta, którego zamierzamy obliczyć
            w dalszej części zadania, \\
    \rho & - & gęstość materiału, \\
    \sigma & - & naprężenie, \\
    \Psi(x,t) & - & średnie przemieszczenie wycinka pręta.
\end{conditions*}


%%%
\subsection{Automata diagram}

\begin{figure}[H]
    \tikzset{
        ->, % makes the edges directed
        >=stealth', % makes the arrow heads bold
        node distance=3cm, % specifies the minimum distance between two nodes. Change if necessary.
        every state/.style={thick, fill=gray!10}, % sets the properties for each ’state’ node
        initial text=$ $, % sets the text that appears on the start arrow
    }
    \centering
    \begin{tikzpicture}
        \node[state, initial] at (-1.763, 2.427)     (q1)      {0000};
        \node[state] at (1.763, 2.427)      (q2)      {0001};
        \node[state] at (2.853, -0.927)     (q3)      {0011};
        \node        at (0, -3)       (q4)  {$\cdots$}; 
        \node[state] at (-2.853, -0.927)    (q5)      {1000};



        \draw   (q1)    edge[bend left, above]  (q2)
                (q2)    edge[bend left, above]  (q3)
                (q3)    edge[bend left, above]  (q4)
                (q4)    edge[bend left, above]  (q5)
                (q5)    edge[bend left, above]  (q1); % A 0 lopp
                % (q1)    edge[bend left, above]  node{1} (q2) % A 1 to B

                % (q2)    edge[bend left, below]  node{0} (q1) % B 0 to A
                % (q2)    edge[right]             node{1} (q3) % B 1 to C

                % (q3)    edge[loop below]        node{1} (q3) % C 1 loop
                % (q3)    edge[above]             node{0} (q4) % C 0 to D

                % (q4)    edge[above]             node{1} (q5) % D 1 to E
                % (q4)    edge[right]             node{0} (q1) % D 0 to A
                
                % (q5)    edge[above]             node{0} (q1); % E 0 to A




    \end{tikzpicture}
    \caption{Schemat automatu}
\end{figure}


\begin{figure}[H]
    \tikzset{
        ->, % makes the edges directed
        >=stealth', % makes the arrow heads bold
        node distance=3cm, % specifies the minimum distance between two nodes. Change if necessary.
        every state/.style={thick, fill=gray!10}, % sets the properties for each ’state’ node
        initial text=$ $, % sets the text that appears on the start arrow
    }
    \centering
    \begin{tikzpicture}
        \node[state, initial] at (0, 0) (q1) {A};
        \node[state] at (4, 0) (q2) {B};
        \node[state, accepting] at (0, -4) (q5) {E};
        \node[state] at (4, -4) (q3) {C};
        \node[state] at (2, -2) (q4) {D};

        \draw   (q1)    edge[loop above]        node{0} (q1) % A 0 lopp
                (q1)    edge[bend left, above]  node{1} (q2) % A 1 to B

                (q2)    edge[bend left, below]  node{0} (q1) % B 0 to A
                (q2)    edge[right]             node{1} (q3) % B 1 to C

                (q3)    edge[loop below]        node{1} (q3) % C 1 loop
                (q3)    edge[above]             node{0} (q4) % C 0 to D

                (q4)    edge[below]             node{1} (q5) % D 1 to E
                (q4)    edge[bend left, below]  node{0} (q1) % D 0 to A
                
                (q5)    edge[left]              node{0} (q1) % E 0 to A
                (q5)    edge[below]             node{1} (q3); % E 1 to C




    \end{tikzpicture}
    \caption{Schemat automatu}
\end{figure}


%%%
\subsection{Karnaugh map}

\begin{figure}[H]
    \centering
    \begin{Karnaugh}{$Q_0 X$}{$Q_2 Q_1$}
        \contingut{0,1,0,0,1,0,0,0,0,0,-,-,-,-,-,-}
        \implicant{1}{1}{green}
        \implicant{4}{12}{red}
    \end{Karnaugh}
    \caption{Tabela Karnaugh dla bitu 0}
\end{figure}

\begin{figure}[H]
    \centering
    \begin{Karnaugh}{$Q_0 X$}{$Q_2 Q_1$}
        \contingut{0,0,0,1,1,1,0,0,0,1,-,-,-,-,-,-}
        \implicant{3}{3}{green}
        \implicant{4}{13}{red}
        \implicant{13}{11}{blue}
    \end{Karnaugh}
    \caption{Tabela Karnaugh dla bitu 1}
\end{figure}

\begin{figure}[H]
    \centering
    \begin{Karnaugh}{$Q_0 X$}{$Q_2 Q_1$}
        \contingut{0,0,0,0,0,0,0,1,0,0,-,-,-,-,-,-}
        \implicant{7}{15}{green}
    \end{Karnaugh}
    \caption{Tabela Karnaugh dla bitu 2}
\end{figure}

\end{document}
